\clearpage
\SourceLeaf{189}
\pdfbookmark[1]{Arrangement and use of the tables.}{witwen-einrichtung}
\begin{center}ARRANGEMENT AND USE OF THE TABLES.\end{center}
\noindent\hspace*{1.2em} [The numerical data in these tables are based on the observations compiled by \textsc{Brune} in volume 16 of \textsc{Crelle}'s Journal für Mathematik concerning the successive admission, from 1776 to 1834, of 31500 married couples to the Royal Prussian General Widows' Pension Institution. The tables give the logarithms of the numbers of women (\(\log fm\)) and men (\(\log FM\)) who, out of 10000 who had completed their 20\textsuperscript{th} year, survived to the end of the year of age (\(m\) or \(M\)) indicated in the middle. There is, however, this departure from \textsc{Brune}: the deaths of men over 80 are calculated in the same proportion as those observations give for women, because, as mentioned in the balance calculation of 1845, the records of the Prussian widows' fund contain insufficient data to determine directly the mortality of men at advanced ages. Beside the logarithms of the living, under the heading decr., stand the absolute values of the differences of those logarithms (\(\log gm=\log\frac{f(m-1)}{fm}\) and \(\log Gm=\log\frac{F(M-1)}{FM}\)), expressed in units of the seventh decimal place. Using the mortality table thus obtained, \textsc{Gauss} calculated the single life and joint life annuities at interest rates of \(3\tfrac{1}{2}\) per cent and 4 per cent.

The single life annuity values, given both as logarithms and as numbers (\(\varphi m,\Phi M\)), take the end of the year of age of the woman (\(m\)) or man (\(M\)) indicated in the middle as the present moment, and assume that one monetary unit is paid at the end of each subsequent year if the person survives to that time, so that
\[
\begin{aligned}
\varphi m\mathord{.}fm&=\rho f(m+1)+\rho\rho f(m+2)+\rho^3f(m+3)+.\ .\ .\\
\Phi M\mathord{.}FM&=\rho F(M+1)+\rho\rho F(M+2)+\rho^3F(M+3)+.\ .\ .
\end{aligned}
\]
where \(\rho\) denotes the discount factor (\(=\tfrac{200}{207}\) at \(3\tfrac{1}{2}\) per cent and \(=\tfrac{25}{26}\) at 4 per cent).

The joint life annuity table contains the logarithms of the values \(\psi(m,M)\). At the moment corresponding to the man's age (\(M\)) given at the side, these are the values of annuities paying one monetary unit at the end of each following year if the man (present age \(=M\)) and woman (age difference \(=m-M\)) are both alive. They are therefore determined by the formula:
\[
\begin{aligned}
\psi(m,M)\mathord{.}fm\mathord{.}Fm&=\rho f(m+1)\mathord{.}F(M+1)\\
&\quad+\rho\rho f(m+2)\mathord{.}F(M+2)\\
&\quad+\rho^3f(m+3)\mathord{.}F(M+3)+.\ .
\end{aligned}
\]
\SourceLeaf{190}
\Needspace{5\baselineskip}
\pdfbookmark[2]{Values of life annuities and life insurance for men.}{witwen-maenner}
\begin{center}\textit{Values of life annuities and life insurance for men.}\end{center}
\noindent\hspace*{1.2em} \textsc{Gauss} used the table of men's life annuities to calculate precisely the effect that the provision in the statutes prohibiting a living member from withdrawing would have. He finds that, for the 42 married members in the balance calculation of 1 Oct. 1845, it would increase the present value of the contributions by 898 thalers at \(3\tfrac{1}{2}\) per cent and by 665 thalers at 4 per cent. He also used this table for some calculations of life insurance values, taking, for a man's present age of \(M\) years, \(\rho-(1-\rho)\Phi M\) as the present value of one monetary unit payable at the end of the year of death.]

\Needspace{6\baselineskip}
\pdfbookmark[2]{Value of the existing widow's pension Pm.}{witwen-pension-pm}
\begin{center}\textit{Value of the existing widow's pension \(Pm\).}\end{center}
\noindent \(m\)\quad widow's present age.\par
\noindent \(\varphi m\)\quad value of the pension if payment is made annually and only to those still living.\par
\noindent \(\rho\)\quad discount factor (\(=\tfrac{25}{26}\) for 4 per cent, \(=\tfrac{200}{207}\) for \(3\tfrac{1}{2}\) per cent), or the interest rate understood in the sense that \(\rho\) grows to 1 after one year.\par
\noindent \(fx\)\quad persons living at age \(x\), as given in the mortality table.

[The balance is calculated for 1 October of a specified year, and this moment is everywhere called simply the present moment. Under the Regulations of 11 October 1833 and the subsequent orders concerning the professors' widows' fund, the pension is paid in half-yearly instalments due on 1 April and 1 October of each year, and for widows it ceases with the month of death, which is paid in full. Accordingly, for the probable present value \(Pm\) of a widow's pension paying one monetary unit annually:]
\[
\begin{aligned}
fm\mathord{.}Pm&=\tfrac{1}{12}\rho^{\frac{1}{2}}\{fm+f(m+\tfrac{1}{12})+f(m+\tfrac{2}{12})+.\ .+f(m+\tfrac{5}{12})\}\\
&\quad+\tfrac{1}{12}\rho\{f(m+\tfrac{6}{12})+f(m+\tfrac{7}{12})+.\ .\ .\ .+f(m+\tfrac{11}{12})\}\\
&\quad+\text{etc.}
\end{aligned}
\]
\noindent or
\[
fm\mathord{.}Pm=\tfrac{1}{2}\rho^{\frac{1}{2}}f(m+\tfrac{5}{24})+\tfrac{1}{2}\rho f(m+\tfrac{17}{24})+\tfrac{1}{2}\rho^{\frac{3}{2}}f(m+\tfrac{29}{24})+\text{etc.}
\]
\noindent for which one may take
\[
=\rho^{\frac{3}{4}}f(m+\tfrac{11}{24})+\rho^{\frac{7}{4}}f(m+\tfrac{35}{24})+\text{etc.}
\]
\noindent The table gives
\[
fm\mathord{.}\varphi m=\rho f(m+1)+\rho\rho f(m+2)+.\ .\ .
\]
\noindent thus
\[
\begin{aligned}
Pm&=\rho^{-\frac{1}{4}}\mathord{.}\frac{f(m-\tfrac{13}{24})}{fm}\mathord{.}\varphi(m-\tfrac{13}{24})\\
&=\rho^{-\frac{1}{4}}\mathord{.}g(m+\tfrac{11}{48})^{\frac{13}{24}}\mathord{.}\varphi(m-\tfrac{13}{24})
\end{aligned}
\]
\relax
[if the values of \(\log f(m-1)-\log fm\) given in the above table under decr. in units of the seventh decimal place are denoted by \(\log\mathord{.}gm\). A formula somewhat more convenient for calculating auxiliary tables is obtained if, for every integer \(n\) and every proper fraction \(k\), \(f(n+k)\) is replaced by the quantity \((k-1)fn+kf(n+1)\), namely:]
\SourceLeaf{191}
\[
Pm=A+B\mathord{.}\varphi m
\]
\noindent where
\[
48A=19\rho^{\frac{1}{2}}+7\rho,\qquad48B=5\rho^{-\frac{1}{2}}+17+19\rho^{\frac{1}{2}}+7\rho
\]
\noindent and, with sufficient approximation
\[
A=\tfrac{13}{24}\rho^{\frac{33}{52}},\qquad B=\rho^{\frac{7}{24}}
\]
\noindent for \(3\tfrac{1}{2}\) per cent.\quad \(\log\tfrac{1}{\rho}=0.0149403\)\quad Exact
\[
A=0.52998474,\qquad\log B=9.9956901
\]
\noindent the approximate formula gives
\[
A=0.5299694,\qquad\log B=9.9956424
\]
\noindent The calculation had used \(\tfrac{13}{24}+\varphi m=P^0\)\par
\noindent one may therefore put
\[
\begin{aligned}
P&=\rho^{\frac{7}{24}}P^0-\tfrac{13}{24}(\rho^{\frac{7}{24}}-\rho^{\frac{3}{4}})\\
&=P^0-\tfrac{7}{24}(1-\rho)P^0-\tfrac{143}{576}(1-\rho)
\end{aligned}
\]
\noindent thus for \(3\tfrac{1}{2}\) per cent.\quad \(P=P^0-\tfrac{49}{4968}P^0-\tfrac{1091}{119232}\)\par
\noindent for 4 per cent.\quad \(P=P^0-\tfrac{7}{624}P^0-\tfrac{11}{1152}\)

\Needspace{14\baselineskip}
Example: no. 65. H.\ . 4 per cent.\quad \(m=53.784\)
\begin{center}
\fontsize{8}{13}\selectfont\setlength{\tabcolsep}{2pt}
\begin{tabular}{ll}
\(\log\varphi(m-\tfrac{13}{24})\ .\ .\ .\ .\ .\)&1.05488\\
\(\text{comp.}\log\rho^{\frac{1}{4}}\ .\ .\ .\ .\ .\ .\)&0.00426\\
\(\log g(m+\tfrac{11}{49})^{\frac{13}{28}}\ .\ .\ .\ .\)&0.00440\\\cline{2-2}
&1.06354\\
\(P=\)&11.5755
\end{tabular}\qquad
\begin{tabular}{ll}
\(\log\varphi m\ .\ .\ .\ .\)&1.04830\\
\(\log\rho^{\frac{7}{24}}\ .\ .\)&\(-0.00497\)\\
\(B\varphi m\ .\ .\ .\ .\)&11.0492\\
\(\tfrac{13}{24}\rho^{\frac{3}{4}}\ .\ .\ .\ .\ .\)&0.5260\\\cline{2-2}
&11.5752
\end{tabular}\par\medskip
Correction of the earlier value \(P^0\)\par\smallskip
\(P^0-\tfrac{1}{99}\tfrac{105}{104}P^0-\tfrac{11}{1152}\)\par\smallskip
\begin{tabular}{r}
\(\tfrac{13}{24}+\varphi m=11.7179\)\\
\(-0.1314\ 5\)\\
\(-0.0095\ 5\)\\\hline
11.5769
\end{tabular}
\end{center}

\Needspace{5\baselineskip}
\pdfbookmark[2]{Existing marriage.}{witwen-stehende-ehe}
\begin{center}\textit{Existing marriage.}\end{center}
\noindent\hspace*{1.2em} The man's age is \(M\) and the woman's \(m\) as of Oct. 1. In the mortality table, living men of age \(x=Fx\); value of the widow's pension \(=R-Q\).

[Under the statutes, the pension begins when the quarter of grace ends, is paid in half-yearly instalments due on 1 April and 1 October each year, and ceases at the end of the month of death. Thus, for the present value \(R-Q\) of a widow's pension that may become payable, amounting to one monetary unit annually,
\[
\begin{aligned}
(R-Q)\mathord{.}fm\mathord{.}Fm&=\tfrac{1}{12}\rho\{Fm-F(m+\tfrac{1}{4})\}\\
&\qquad\{f(m+\tfrac{6}{12})+f(m+\tfrac{7}{12})+f(m+\tfrac{8}{12})\}\\
&\quad+\tfrac{1}{12}\rho\{Fm-F(m+\tfrac{2}{4})\}\\
&\qquad\{f(m+\tfrac{9}{12})+f(m+\tfrac{10}{12})+f(m+\tfrac{11}{12})\}\\
&\quad+\tfrac{1}{12}\rho^{\frac{3}{2}}\{Fm-F(m+\tfrac{3}{4})\}\\
&\qquad\{f(m+\tfrac{12}{12})+f(m+\tfrac{13}{12})+f(m+\tfrac{14}{12})\}\\
&\quad+\tfrac{1}{12}\rho^{\frac{3}{2}}\{Fm-F(m+\tfrac{4}{4})\}\\
&\qquad\{f(m+\tfrac{15}{12})+f(m+\tfrac{16}{12})+f(m+\tfrac{17}{12})\}\\
&\quad+\text{etc.}
\end{aligned}
\]
\noindent or:]
\[
\begin{aligned}
R\mathord{.}fm&=\tfrac{1}{2}\rho f(m+\tfrac{17}{24})+\tfrac{1}{2}\rho^{\frac{3}{2}}f(m+\tfrac{29}{24})+\tfrac{1}{2}\rho\rho f(m+\tfrac{41}{24})+.\ .\\
&=\rho^{\frac{5}{4}}f(m+\tfrac{23}{24})+\rho^{\frac{9}{4}}f(m+\tfrac{47}{24})+.\ .
\end{aligned}
\]
\[
R=\tfrac{1}{24}\rho^{\frac{5}{4}}+\rho^{\frac{7}{24}}\varphi m
\]
\[
\begin{aligned}
Q\mathord{.}fm\mathord{.}Fm&=\rho^{\frac{5}{4}}f(m+\tfrac{23}{24})\mathord{.}F(M+\tfrac{5}{8})\\
&\quad+\rho^{\frac{9}{4}}f(m+\tfrac{47}{24})F(M+\tfrac{13}{8})+.\ .
\end{aligned}
\]
\SourceLeaf{192}
\Needspace{5\baselineskip}
\pdfbookmark[2]{Value of the contributions still to be paid S.}{witwen-beitraege}
\begin{center}\textit{Value of the contributions still to be paid \(S\).}\end{center}
\noindent\hspace*{1.2em} [Regulations of 11 October 1833 and 24 November 1846. Every participant in the widows' fund must pay the contribution to the accountant annually in arrears on 17 September, and these contributions are reckoned from Michaelmas to Michaelmas. — If a member of the widows' fund dies in the first half of the contribution year, that is, in the months from October to March inclusive, the heirs no longer have to pay the contribution for that year; but if a member dies in the second half of the year, the heirs must still pay the full contribution on 17 September of the year of death. — Participants must give notice of withdrawal by a written declaration before 17 September of the membership year running from Michaelmas to Michaelmas; the annual contribution falling due on that day must nevertheless be paid once more in full. Anyone who misses this deadline must still pay for the whole following membership year, but also remains a member until then.

For the present value \(S\) of one monetary unit payable annually as a contribution, one therefore obtains:]
\[
S\mathord{.}fm\mathord{.}Fm=\rho\mathord{.}fm\mathord{.}F(M+\tfrac{1}{2})+\rho\rho\mathord{.}f(m+1)\mathord{.}F(M+\tfrac{3}{2})+.\ .
\]
\noindent Thus, if one writes
\[
\psi(m,M)\mathord{.}fm\mathord{.}Fm=\rho f(m+1)F(m+1)+\rho\rho f(m+2)F(m+2)+.\ .\ .
\]
\noindent then
\[
\begin{aligned}
Q&=\frac{f(m-\tfrac{1}{24})}{fm}\mathord{.}\frac{F(M-\tfrac{3}{8})}{FM}\mathord{.}\rho^{\frac{1}{4}}\psi(m-\tfrac{1}{24},M-\tfrac{3}{8})\\[5pt]
S&=\frac{f(m-1)}{fm}\mathord{.}\frac{F(M-\tfrac{1}{2})}{FM}\mathord{.}\psi(m-1,M-\tfrac{1}{2})
\end{aligned}
\]
\noindent If one also writes
\[
\begin{aligned}
\frac{f(m-1)}{fm}&=gm\\[5pt]
\frac{F(M-1)}{FM}&=GM
\end{aligned}
\]
\noindent one obtains
\[
\begin{aligned}
R&=\tfrac{1}{24}\rho^{\frac{5}{4}}+\rho^{\frac{7}{24}}\varphi m\\
Q&=g(m+\tfrac{23}{48})^{\frac{1}{24}}\mathord{.}G(M+\tfrac{5}{16})^{\frac{3}{8}}\mathord{.}\rho^{\frac{1}{4}}\mathord{.}\psi(m-\tfrac{1}{24},M-\tfrac{3}{8})\\
S&=gm\mathord{.}G(M+\tfrac{1}{4})^{\frac{1}{2}}\mathord{.}\psi(m-1,M-\tfrac{1}{2})
\end{aligned}
\]
\SourceLeaf{193}
\Needspace{17\baselineskip}
Example: no. 138 Conradi.\quad \(M=65.02,\quad m=46.08\)
\begin{center}
\fontsize{9}{14}\selectfont\setlength{\tabcolsep}{5pt}
\begin{tabular}{lrr}
&4 per cent.&\(3\tfrac{1}{2}\) per cent.\\
\(\log\varphi m\ .\ .\ .\ .\)&1.12875&1.15233\\
\(\tfrac{7}{24}\log\rho\ .\ .\ .\ .\)&\(-497\)&\(-436\)\\
&13.2979&14.0594\\
\(\tfrac{1}{24}\rho^{\frac{5}{4}}\ .\ .\ .\ .\)&0.0397&399\\\cline{2-3}
\(R=\)&13.3376&14.0993\\\cline{2-3}
for \(Q,\ \log\psi\ .\ .\ .\)&0.80912&0.82183\\
\(\log gG\ .\ .\)&\(+897\)&\(+897\)\\
\(\tfrac{1}{4}\log\rho\ .\ .\ .\)&\(-426\)&\(-374\)\\
\(Q=\)&6.5137&6.7150\\
\(R-Q=\)&6.8239&7.3843\\\hline
for \(S,\ \log\psi\ .\ .\ .\)&0.81325&0.82606\\
\(\log gG\ .\ .\)&\(+1744\)&\(+1744\)\\
\(S=\)&6.7716&6.9743\\\hline
\end{tabular}
\end{center}
\[
\begin{aligned}
\log\rho&=-0.0170333393\quad\text{for 4 per cent.}\\
\log\rho&=-0.0149403498\quad\text{for }3\tfrac{1}{2}\text{ per cent.}
\end{aligned}
\]
\begin{center}
\fontsize{9}{14}\selectfont\setlength{\tabcolsep}{4pt}
\begin{tabular}{c|c}
for \(Q\)&for \(S\)\\
\(\psi(46.04;64.65)\ .\ .\ 18.61\)&\(\psi(45.08;64.52)\ .\ .\ 19.44\)\\
18.61&19.44\\
\begin{tabular}{r||r|r}64&0.82079&0.83374\\65&0.80284&0.81540\end{tabular}&\begin{tabular}{r|r}0.82251&0.83552\\0.80470&0.81732\end{tabular}\\
\(g(46.56)\ .\ .\ .\ 25\)&\(g(46.08)\ .\ .\ .\ 588\)\\
\(G(65.33)\ .\ .\ .\ 872\)&\(G(65.27)\ .\ .\ .\ 1156\)
\end{tabular}
\end{center}

\noindent\hspace*{1.2em} In constructing these tables, \textsc{Gauss} used his five-place table for calculating the logarithm of the sum of quantities given only by their logarithms rather than by the quantities themselves, and therefore began, for each interest rate and each age difference, by determining the annuity values for the highest ages. For this reason, even apart from the improvement that the mortality table has already undergone since then through the broader observations of the Prussian widows' fund, individual numerical entries may be inaccurate by a few units in the fifth decimal place of the logarithms.

I have corrected in printing some computational errors that I noticed by forming the quotients between the \(3\tfrac{1}{2}\) and 4 per cent annuity values and the differences between successive quotients. In units of the fifth decimal place of the logarithm, these errors amount, at the places where they originate, to: \(+20\) for the woman's 77th year in her single life annuity values at 4 per cent, and also for the man's 71st and 93rd years and the woman's respective age differences of \(+1\) and 0 years at \(3\tfrac{1}{2}\) per cent; \(+10\) for the man's 76th year in his single life annuity value at \(3\tfrac{1}{2}\) per cent, and the same amount for the man's 79th year and the woman's age difference of \(-9\) years at 4 per cent; \(-10\) for the man's 28th, 68th, 91st, 94th and 94th years and the respective age differences of \(-8\), \(-13\), \(-11\), \(-2\), and \(-14\) years at the respective rates \(4.4.3\tfrac{1}{2}.4\) and \(3\tfrac{1}{2}\) per cent. These and a few other smaller computational errors have had some influence on the determination of the annuity values for the immediately younger ages; this influence disappears gradually and, in the extreme case, only at the end of the second decade. The numerical values of the single life annuities are derived from the logarithms, and after the stated correction of the latter they have also had to undergo a corresponding change here.

The formulas to be used in applying the tables, and their worked examples, are taken from scattered notes on individual slips and are explained here by a few inserted passages
\begin{flushright}\textsc{Schering.}\end{flushright}
\begin{center}\rule{.3\linewidth}{.3pt}\end{center}
